\documentclass[9pt,aspectratio=169,xcolor=table]{beamer}

\usepackage[sfdefault]{roboto}
\usepackage[utf8]{inputenc}
\usepackage[T1]{fontenc}

\usepackage{styles/fluxmacros}
\usefolder{styles}
\usetheme[style=red]{flux}

\definecolor{primaryLight}{HTML}{9B2242}
\definecolor{primary}{HTML}{7B1E3A}
\definecolor{secondary}{HTML}{3A5A6B}
\definecolor{tertiary}{HTML}{8C6A3F}

\usepackage{booktabs}
\usepackage{colortbl}
\usepackage{ragged2e}
\usepackage{amssymb}
\usepackage{multirow}
\usepackage{tikz}
\usetikzlibrary{arrows.meta, positioning, fit, backgrounds, calc, shapes.geometric}
\setbeamertemplate{itemize items}[square]
\usepackage[scaled=.92]{beramono}
\usepackage{listings}
\usepackage{array}
\definecolor{stripe}{gray}{0.94}
\newcommand{\striped}{\rowcolors{1}{white}{stripe}}

\definecolor{codebg}{HTML}{FBF2F4}
\definecolor{codegreen}{rgb}{0,0.45,0}
\definecolor{codestring}{rgb}{0.58,0.1,0.2}
\lstdefinestyle{skpython}{
    basicstyle=\ttfamily\footnotesize,
    keywordstyle=\color{primary}\bfseries,
    commentstyle=\itshape\color{codegreen},
    stringstyle=\ttfamily\color{codestring},
    backgroundcolor=\color{codebg},
    breaklines=true,
    keepspaces=true,
    showspaces=false,
    showstringspaces=false,
    frame=single,
    rulecolor=\color{primary!40},
    numbers=none,
}
\lstdefinestyle{skoutput}{
    basicstyle=\ttfamily\footnotesize,
    backgroundcolor=\color{secondary!8},
    breaklines=true,
    keepspaces=true,
    showspaces=false,
    showstringspaces=false,
    frame=single,
    rulecolor=\color{secondary!50},
    numbers=none,
}
\lstset{language=Python, style=skpython}

\definecolor{cycfill}{HTML}{F3DCE4}
\definecolor{cycline}{HTML}{7B1E3A}
\definecolor{cycdofill}{HTML}{E8EEF0}
\definecolor{cycdoline}{HTML}{3A5A6B}
\tikzset{
  cycstep/.style = {font=\sffamily\small, align=center, rounded corners=2pt,
                     draw=cycline, fill=cycfill, line width=0.7pt,
                     minimum height=7mm, minimum width=20mm, inner sep=3pt},
  cycdo/.style   = {cycstep, fill=cycdofill, draw=cycdoline},
  flow/.style    = {-{Stealth[length=2mm]}, thick, draw=cycline},
}

\AtBeginSection[]
{
{
    \setbeamercolor{background canvas}{bg=primaryLight!6}
    \begin{frame}[plain]
        \frametitle{Table of Contents}
        \tableofcontents[currentsection]
    \end{frame}
}
}

\title{Decisions and Repetition}
\subtitle{\emph{Engineering Computation with Python} --- Chapter 4}
\author{Mun\'{i}m Zabidi}
\institute{\color{primary}Faculty of Electrical Engineering, Universiti Teknologi Malaysia}
\date{\today}
\titlegraphic{assets/logoutmwhite.png}

\begin{document}

\titlepage

\begin{frame}{Table of Contents}
    \tableofcontents
\end{frame}

%===============================================================================
\section{Boolean Foundations}
%===============================================================================

\begin{frame}{From Straight-Line Code to Real Decisions}
    \leavevmode
    Every program so far has run straight through, one line at a time.
    That is not enough for engineering work, which is full of:

    \vspace{2mm}
    \begin{columns}[t]
        \begin{column}{.5\textwidth}
            \begin{block}{Decisions}
                \small Is this reading within tolerance?
            \end{block}
        \end{column}
        \begin{column}{.5\textwidth}
            \begin{block}{Repetition}
                \small Apply this check to all $101$ samples.
            \end{block}
        \end{column}
    \end{columns}

    \vspace{3mm}
    \begin{alertblock}{The foundation beneath both}
        Every decision and every loop condition is built from
        \alert{Boolean expressions} --- expressions that evaluate to
        \texttt{True} or \texttt{False}.
    \end{alertblock}
\end{frame}

\begin{frame}{Relational and Logical Operators}
    \leavevmode
    \small
    \begin{columns}[t]
        \begin{column}{.45\textwidth}
            \striped
            \begin{center}
            \begin{tabular}{cl}
                \toprule
                \textbf{Symbol} & \textbf{Meaning} \\
                \midrule
                \texttt{<}  & Less than \\
                \texttt{<=} & Less than or equal \\
                \texttt{==} & Equal \\
                \texttt{!=} & Not equal \\
                \texttt{>}  & Greater than \\
                \texttt{>=} & Greater than or equal \\
                \bottomrule
            \end{tabular}
            \end{center}
        \end{column}
        \begin{column}{.5\textwidth}
            \striped
            \begin{center}
            \begin{tabular}{lcc}
                \toprule
                \textbf{Logical} & \textbf{Elem.-wise} & \textbf{Short-circ.} \\
                \midrule
                AND & \texttt{\&} & \texttt{and} \\
                OR  & \texttt{|}  & \texttt{or}  \\
                NOT & \texttt{\~{}} & \texttt{not} \\
                \bottomrule
            \end{tabular}
            \end{center}
        \end{column}
    \end{columns}

    \vspace{3mm}
    {\small Any nonzero value is \texttt{True}; zero is \texttt{False}.}
\end{frame}

\begin{frame}[fragile]{Logical vs. Bitwise: Don't Confuse Them}
    \leavevmode
    \alert{Logical} operators (\texttt{and}, \texttt{or}, \texttt{not})
    output \texttt{True}/\texttt{False}. \alert{Bitwise} operators
    (\texttt{\&}, \texttt{|}, \texttt{\textasciitilde}) work on the
    binary representation of numbers.

\begin{lstlisting}
True and False  # False
True or False   # True
not True        # False

6 & 3  # 2  (Binary: 110 & 011 = 010)
6 | 3  # 7  (Binary: 110 | 011 = 111)
~5     # -6 (two's complement)
\end{lstlisting}

    \vspace{1mm}
    \begin{alertblock}{A common mistake}
        Using \texttt{\&} where you meant \texttt{and} compiles fine and
        gives a misleading numeric answer instead of an error.
    \end{alertblock}
\end{frame}

\begin{frame}[fragile]{Relational Operators in Action}
    \leavevmode
\begin{lstlisting}
a = 5
b = 10

print(a < b)     # True
print(a <= b)    # True
print(a == b)    # False
print(a != b)    # True
print(a > b)     # False
print(a >= b)    # False
\end{lstlisting}

    \vspace{1mm}
    {\small Each of the six operators applied to the same pair of values
    --- compare the results directly.}
\end{frame}

\begin{frame}{Operator Precedence}
    \leavevmode
    \small
    Python follows a strict precedence order across arithmetic, bitwise,
    and Boolean operators --- \alert{from highest to lowest} below:

    \vspace{2mm}
    \striped
    \begin{center}
    \begin{tabular}{p{4.0cm}p{6.5cm}}
        \toprule
        \textbf{Operator} & \textbf{Description} \\
        \midrule
        \texttt{**} & Exponentiation \\
        \texttt{*}, \texttt{/}, \texttt{//}, \texttt{\%} & Mult., div., floor div., remainder \\
        \texttt{+}, \texttt{-} & Addition, subtraction \\
        \texttt{\&}, \texttt{\^{}}, \texttt{|} & Bitwise AND, XOR, OR \\
        \texttt{<}, \texttt{<=}, \texttt{>}, \texttt{>=}, \texttt{!=}, \texttt{==} & Comparisons \\
        \texttt{not}, \texttt{and}, \texttt{or} & Boolean logic \\
        \bottomrule
    \end{tabular}
    \end{center}

    \vspace{2mm}
    {\small \alert{When in doubt, parenthesise.}}
\end{frame}

\begin{frame}[fragile]{Precedence in Practice}
    \leavevmode
\begin{lstlisting}
a, b, c = 5, 2, 3

expr1 = (a + b) * c        # Parens first: 21
expr2 = a ** b + c         # Power first: 28
expr6 = a < b and b > c    # Relational, then AND: False
expr7 = a | b and c        # Bitwise OR then AND: 3
\end{lstlisting}

    \vspace{2mm}
    \begin{alertblock}{Expression 7 deserves a closer look}
        \texttt{a | b} evaluates first (bitwise OR: $5|2=7$). Since
        \texttt{and} returns one of its \alert{actual operands}, not
        just \texttt{True}/\texttt{False}, the result is \texttt{c = 3}.
    \end{alertblock}
\end{frame}

\begin{frame}[fragile]{Combining All Three Kinds at Once}
    \leavevmode
    Real conditions usually combine more than one operator kind:

\begin{lstlisting}
x, y, z = 8, 5, 12

result = ((x + y) > z) and ((x - y) * 2 < z) or (z % y == 0)
print(result)  # True

# (x + y) > z       : is the sum greater than z?
# (x - y) * 2 < z   : is double the difference less than z?
# z % y == 0        : does z divide evenly by y?
\end{lstlisting}

    \vspace{1mm}
    {\small This is where precedence starts to matter --- work through
    it piece by piece before trusting the printed result.}
\end{frame}

\begin{frame}[fragile]{Chained Comparisons: Readable Range Checks}
    \leavevmode
    A range check written the traditional way:

\begin{lstlisting}
x = 10
if x > 5 and x < 15:
    print("x is between 5 and 15")
\end{lstlisting}

    can be written as a single \alert{chained comparison}:

\begin{lstlisting}
if 5 < x < 15:
    print("x is between 5 and 15")
\end{lstlisting}

    \vspace{1mm}
    {\small Python evaluates left to right and combines the result with
    an implicit \texttt{and} --- exactly the kind of test that decides
    whether a measurement lies within tolerance.}
\end{frame}

\begin{frame}[fragile]{Chained Comparisons: More Examples}
    \leavevmode
\begin{lstlisting}
x = 15
print(10 < x < 20)          # True

y = 25
print(20 < y <= 30)         # True

a = 4
print(2 < a**2 < 20)        # True, 16 is between 2 and 20

name = "Bob"
print("Alice" < name < "Charlie")  # True (alphabetical)
\end{lstlisting}

    \vspace{1mm}
    {\small Chained comparisons can involve more than two operators, mix
    data types, and even call functions --- the chain is still evaluated
    left to right.}
\end{frame}

%===============================================================================
\section{Decisions}
%===============================================================================

\begin{frame}{Three Kinds of Statement}
    \leavevmode
    Program statements fall into three categories:

    \vspace{2mm}
    \begin{enumerate}
        \item \textbf{Assignment} --- \texttt{variable = expression}
        \item \textbf{Decision} --- executes statements if a condition
              is \texttt{True} (\texttt{if}, \texttt{if-else},
              \texttt{if-elif-else})
        \item \textbf{Repetition} --- executes statements many times
              (\texttt{while}, \texttt{for})
    \end{enumerate}

    \vspace{3mm}
    {\small This section covers decisions; the next covers repetition.}
\end{frame}

\begin{frame}[fragile]{The if-elif-else Structure}
    \leavevmode
\begin{lstlisting}
if condition_1:
    statements
elif condition_2:
    statements
else:
    statements
\end{lstlisting}

    \vspace{2mm}
    \begin{block}{Example}
\begin{lstlisting}
x = 2
if x < 0:
    print('neg')
elif x == 0:
    print('zero')
else:
    print('non-neg')   # prints: non-neg
\end{lstlisting}
    \end{block}
\end{frame}

\begin{frame}[fragile]{Multiple Conditions with and}
    \leavevmode
    Logical operators combine multiple conditions in a single
    \texttt{if}. For $ax^2+bx+c=0$, equal roots exist when
    $b^2-4ac=0$ \alert{and} $a \neq 0$:

\begin{lstlisting}
a = float(input('Input a value: '))
b = float(input('Input b value: '))
c = float(input('Input c value: '))

d = b**2 - 4 * a * c

if d == 0 and a != 0:
    x = -b / (2 * a)
    print(f'The value of x is: {x}')
else:
    print('Different root')
\end{lstlisting}
\end{frame}

\begin{frame}[fragile]{Multiple Conditions with if-elif-else}
    \leavevmode
    Three outcomes: $d>0$ (two real roots), $d=0$ (one root), $d<0$
    (complex roots) --- \alert{more than two conditions needs} \texttt{elif}:

\begin{lstlisting}
d = b**2 - 4 * a * c

if d == 0 and a != 0:
    x1 = -b / (2 * a)
    print(f'The roots are both: {x1}')
elif d > 0 and a != 0:
    x1 = (-b + (d**0.5)) / (2 * a)
    x2 = (-b - (d**0.5)) / (2 * a)
    print(f'The roots are: {x1} and {x2}')
else:
    print('Complex root')
\end{lstlisting}
\end{frame}

\begin{frame}[fragile]{Nested if: When One Check Depends on Another}
    \leavevmode
    Check $a \neq 0$ \alert{first}; only then does the discriminant
    calculation make sense (if $a=0$, the equation is linear, not
    quadratic):

\begin{lstlisting}
if a != 0:
    if d == 0:
        x1 = -b / (2 * a)
        print(f'The roots are both: {x1}')
    elif d > 0:
        x1 = (-b + (d**0.5)) / (2 * a)
        x2 = (-b - (d**0.5)) / (2 * a)
        print(f'The roots are: {x1} and {x2}')
    else:
        print('Complex root')
else:
    x1 = -c / b
    print(f'The root is: {x1}')
\end{lstlisting}
\end{frame}

%===============================================================================
\section{Loops}
%===============================================================================

\begin{frame}[fragile]{Two Loop Constructs}
    \leavevmode
    \begin{columns}[t]
        \begin{column}{.5\textwidth}
            \begin{block}{\texttt{for}: determinate}
                \small Repeats for every value in a sequence (e.g.
                \texttt{range}) --- the number of iterations is known
                \alert{before} the loop starts.
\begin{lstlisting}
for i in range(start, stop, step):
    statements
\end{lstlisting}
            \end{block}
        \end{column}
        \begin{column}{.5\textwidth}
            \begin{block}{\texttt{while}: indeterminate}
                \small Repeats as long as a condition is true --- the
                number of iterations is \alert{not} necessarily known
                in advance.
\begin{lstlisting}
while condition:
    statements
\end{lstlisting}
            \end{block}
        \end{column}
    \end{columns}

    \vspace{2mm}
    {\small \texttt{range(j, k)}: generates $j$ to $k-1$.
    \texttt{range(j, k, m)}: step $m$.}
\end{frame}

\begin{frame}[fragile]{Worked Example: Compound Interest with for}
    \leavevmode
    \[
    B = a(1+r)^n
    \]
    Evaluate $B$ for $a=\$100$, $r=8\%$, over $n = 2,4,6,8,10$:

\begin{lstlisting}
a = 100   # initial investment
r = 0.08  # annual interest rate

for n in range(2, 11, 2):
    B = a * (1 + r) ** n
    print(f'n = {n}, B = {B:.2f}')
\end{lstlisting}
\begin{lstlisting}[style=skoutput]
n = 2, B = 116.64
n = 4, B = 136.05
...
n = 10, B = 215.89
\end{lstlisting}
\end{frame}

\begin{frame}{Recognizing $\Sigma$ as a Loop}
    \leavevmode
    Your mathematics courses write $\sum_{i=1}^{5} i = 15$. A \texttt{for}
    loop is precisely this idea, written one step at a time:

    \vspace{2mm}
    \striped
    \begin{center}
    \begin{tabular}{p{3.2cm}p{3.2cm}p{3.6cm}}
        \toprule
        \textbf{Notation} & \textbf{Meaning} & \textbf{Python} \\
        \midrule
        (implicit) & Total starts at 0 & \texttt{total\_sum = 0} \\
        $i=1$ (below $\Sigma$) & Start index & \texttt{range(1, ...} \\
        $5$ (above $\Sigma$) & End index (incl.) & \texttt{..., 6)} \\
        summed term & Added each step & \texttt{total\_sum += i} \\
        \bottomrule
    \end{tabular}
    \end{center}
\end{frame}

\begin{frame}[fragile]{The Summation Pattern}
    \leavevmode
\begin{lstlisting}
total_sum = 0
for i in range(1, 6):     # i = 1, 2, 3, 4, 5
    total_sum += i
print(total_sum)          # 15
\end{lstlisting}

    \vspace{2mm}
    \begin{alertblock}{The same skeleton, every time}
        $\sum i^2$ or $\sum 1/i$ use the \alert{exact same loop} ---
        only the expression added to the running total changes.
        Recognizing $\Sigma$ as ``a \texttt{for} loop with a running
        total'' is one of the most useful shortcuts in this course.
    \end{alertblock}
\end{frame}

\begin{frame}[fragile]{The Same Computation with while}
    \leavevmode
\begin{lstlisting}
a = 100; r = 0.08
n = 2

while n <= 10:
    B = a * (1 + r) ** n
    print(f'n = {n}, B = {B:.2f}')
    n += 2   # must increment manually
\end{lstlisting}

    \vspace{2mm}
    \begin{block}{The key difference from for}
        \texttt{while} requires \texttt{n} to be \alert{initialized and
        manually incremented} --- the number of iterations is governed
        entirely by the condition.
    \end{block}
\end{frame}

\begin{frame}[fragile]{Worked Example: Capacitor Charging}
    \leavevmode
    \[
    Q(t) = CV\left(1 - e^{-t/RC}\right)
    \]
    Monitor charge every $0.5\,\text{s}$ until it exceeds $2$ units
    ($V=9$, $R=4$, $C=1$) --- the number of steps is \alert{not known
    in advance}, so \texttt{while} is the natural choice:

\begin{lstlisting}
import math
V, R, C = 9, 4, 1
t, q = 0, 0
Q, T = [], []

while q <= 2:
    q = C * V * (1 - math.exp(-t / (R * C)))
    Q.append(q); T.append(t)
    t += 0.5
\end{lstlisting}
\end{frame}

\begin{frame}[fragile]{Two Ways to Grow a List}
    \leavevmode
    \begin{columns}[t]
        \begin{column}{.5\textwidth}
            \begin{block}{\texttt{.append()}}
\begin{lstlisting}
Q.append(q)
T.append(t)
\end{lstlisting}
            \end{block}
        \end{column}
        \begin{column}{.5\textwidth}
            \begin{block}{\texttt{+=}}
\begin{lstlisting}
Q += [round(q, 4)]
T += [t]
\end{lstlisting}
            \end{block}
        \end{column}
    \end{columns}

    \vspace{3mm}
\begin{lstlisting}[style=skoutput]
Q: [0.0, 1.0575, 1.9908, 2.8144]
T: [0, 0.5, 1.0, 1.5]
\end{lstlisting}

    \vspace{1mm}
    {\small Both produce identical results --- a matter of style.}
\end{frame}

\begin{frame}[fragile]{Nested Loops: Two Values at Once}
    \leavevmode
    $B$ for 3 values of $a$ \alert{and} 5 values of $n$ needs \alert{two}
    loops, one inside the other:

\begin{lstlisting}
a_values = [100, 500, 800]
n_values = [2, 4, 6, 8, 10]
r = 0.09

for a in a_values:        # outer loop
    for n in n_values:    # inner loop
        B = a * (1 + r) ** n
        print(f"a={a}, n={n}, B={B:.2f}")
\end{lstlisting}

    \vspace{1mm}
    {\small The inner loop runs to completion for \emph{every} value of
    the outer loop --- $3 \times 5 = 15$ total iterations.}
\end{frame}

\begin{frame}[fragile]{continue: Skip This Iteration}
    \leavevmode
    \texttt{continue} skips the rest of the current iteration and moves
    to the next one.

\begin{lstlisting}
for n in range(1, 51):
    if n % 7 != 0:
        continue          # skip non-multiples of 7
    print(f"Divisible by 7: {n}")
\end{lstlisting}
\begin{lstlisting}[style=skoutput]
Divisible by 7: 7
Divisible by 7: 14
...
Divisible by 7: 49
\end{lstlisting}

    \vspace{1mm}
    {\small In nested loops, \texttt{continue} only skips the
    \alert{inner} loop it occurs in.}
\end{frame}

\begin{frame}[fragile]{break: Exit the Loop Entirely}
    \leavevmode
    \texttt{break} terminates the loop immediately --- useful for an
    \alert{indeterminate} loop that should stop as soon as a goal is met.

\begin{lstlisting}
import random
x = random.randint(1, 6)
n = 0

while True:               # loop forever...
    guess = int(input('Guess the number: '))
    n += 1
    if guess == x:
        print(f'Correct after {n} tries!')
        break              # ...until the guess is right
\end{lstlisting}
\end{frame}

\begin{frame}{break vs. continue}
    \leavevmode
    \begin{columns}[t]
        \begin{column}{.5\textwidth}
            \begin{block}{break}
                \small Terminates the loop entirely. Statements after
                \texttt{break} inside the loop do not execute.
            \end{block}
        \end{column}
        \begin{column}{.5\textwidth}
            \begin{block}{continue}
                \small Skips to the next iteration. Remaining statements
                in \emph{this} iteration are skipped.
            \end{block}
        \end{column}
    \end{columns}

    \vspace{3mm}
    \begin{alertblock}{A caution}
        Both are useful, but a loop whose logic depends heavily on them
        is often clearer when rewritten without them.
    \end{alertblock}
\end{frame}

%===============================================================================
\section{Summary}
%===============================================================================

\begin{frame}{What You Should Take Away}
    \leavevmode
    \small
    \begin{itemize}
        \item \textbf{Relational operators} return \texttt{True}/\texttt{False};
              these Boolean values are what every decision and loop
              condition is built from.
        \item \textbf{Logical} (\texttt{and}, \texttt{or}, \texttt{not})
              vs. \textbf{bitwise} (\texttt{\&}, \texttt{|},
              \texttt{\textasciitilde}) --- confusing the two is a
              frequent source of error.
        \item Python evaluates by strict \textbf{precedence} ---
              arithmetic, then comparison, then logical. Parenthesise
              when in doubt.
        \item \textbf{Chained comparisons} (\texttt{5 < x < 15}) express
              range checks in one readable line.
    \end{itemize}
\end{frame}

\begin{frame}{What You Should Take Away (continued)}
    \leavevmode
    \small
    \begin{itemize}
        \item \textbf{if-elif-else} selects which block runs; an
              \texttt{elif} chain stops at the first true condition, so
              order matters.
        \item \textbf{for} runs a known number of times over a sequence;
              \textbf{while} runs until a condition is false --- the
              right choice when iterations are not known in advance.
        \item \texttt{break} exits a loop immediately; \texttt{continue}
              skips to the next iteration.
        \item A $\Sigma$ in a formula is ``a \texttt{for} loop with a
              running total'' --- one of the most useful shortcuts in
              this course.
    \end{itemize}

    \vspace{4mm}
    \begin{center}
        \large\color{primary}\textbf{Next: Functions and Reusable Calculations}
    \end{center}
\end{frame}

\end{document}
